\documentclass[11pt]{article}
\usepackage[margin=1in]{geometry}
\usepackage{amsmath,amssymb,amsthm}
\usepackage{mathtools}
\usepackage{booktabs}
\usepackage{hyperref}

\newtheorem{theorem}{Theorem}[section]
\newtheorem{proposition}[theorem]{Proposition}
\newtheorem{lemma}[theorem]{Lemma}
\newtheorem{corollary}[theorem]{Corollary}
\newtheorem{conjecture}[theorem]{Conjecture}
\theoremstyle{definition}
\newtheorem{definition}[theorem]{Definition}
\newtheorem{remark}[theorem]{Remark}

\newcommand{\KL}{\mathrm{KL}}
\newcommand{\Var}{\mathrm{Var}}
\newcommand{\Cov}{\mathrm{Cov}}
\newcommand{\tr}{\mathrm{tr}}
\newcommand{\gap}{\mathrm{gap}}
\newcommand{\Hess}{\mathrm{Hess}}
\newcommand{\su}{\mathfrak{su}}
\newcommand{\R}{\mathbb{R}}
\newcommand{\Z}{\mathbb{Z}}
\newcommand{\E}{\mathbb{E}}

\title{Entropy-controlled renormalization for lattice Yang--Mills: \\
geometric decay of non-Gaussianity at all scales}
\author{Ramiro da Costa Fontes\thanks{Quijotic Research.
Email: \texttt{quijoticresearch@gmail.com}}}
\date{February 2026}

\begin{document}
\maketitle

\begin{abstract}
We introduce an entropy-based framework for tracking non-Gaussianity 
through block-spin renormalization group (RG) transformations in 
lattice Yang--Mills theory.  The central object is the relative entropy 
$W_k = \KL(\mu_k \| \nu_k)$ between the effective Yang--Mills measure 
$\mu_k$ at RG scale~$k$ and its Gaussian approximation $\nu_k$, 
whose precision is the iterated Schur complement of the lattice 
Hessian.  We establish four structural results valid at all couplings: 
(i)~the Schur complement factorizes as $S = B^\top B$ (preserving 
incidence structure under RG), (ii)~the exact comparability 
$R\,S\,R^\top = L^{-d} H_{\mathrm{coarse}}$ (proved via a 
face-averaging cancellation $H_{yx}R^\top = 0$), (iii)~the 
KL chain rule $W_{k+1} = W_k - D_k$ with $D_k \geq 0$, and 
(iv)~a per-plaquette variance bound 
$\Var_{\nu_k}(\phi_p) \leq 1/\beta$ at every scale~$k$, 
following from the projection structure of $B(B^\top B)^+ B^\top$.

The per-plaquette bound ensures the quartic approximation to the 
Wilson action is valid at \emph{all} RG scales (not just a 
perturbative window), despite the per-link variance growing as 
$L^{kd}/\beta$.  Combined with $W_0 = O(g^4)$ and the analytical 
dissipation bound $D_k/W_k \geq 0.527$, this gives geometric 
entropy decay $W_k \leq (1 - f)^k C/\beta^2$ through the entire 
RG flow.

We identify the Gaussian Schur complement as the analogue of 
Perelman's Gaussian shrinker, the KL chain rule as the analogue 
of $W$-entropy monotonicity, and the antisymmetric Christoffel 
cancellation (proved in the companion paper) as the structural 
identity that ensures $W_0 = O(g^4)$ rather than $O(g^2)$.  
For compact $\mathrm{U}(1)$, combining the entropy framework with 
the Villain decomposition (Poisson resummation) recovers the 
correct physics in all dimensions---including the mass gap in $d = 3$ 
via Debye screening of the monopole gas.  For $\mathrm{SU}(N)$ in 
$d = 4$, the mass gap requires a non-abelian analog of the Villain 
decomposition, for which the character expansion is a natural candidate.  
We outline a path to the mass gap via two-regime 
composition of log-Sobolev inequalities (Lu--Yau method): 
at weak coupling, the per-plaquette bound controls the 
Holley--Stroock perturbation of the Gaussian reference; at 
strong coupling, the compact group provides the spectral gap 
via Holley--Stroock perturbation of the Haar measure with no 
Gaussian input.  The block-spin RG (marginalization over interior 
links) is well-defined on the full compact group at all scales; 
the Gaussian serves only as a tracking reference.  Combining 
the two regimes yields a volume-uniform spectral gap 
$\rho = \Omega(1)$ in coarse units and a physical mass gap 
$m \geq c\,\beta^{-1/d}/a$, conditional on verification of 
the Hessian Dobrushin threshold at strong coupling.
\end{abstract}

\tableofcontents

%% ================================================================
\section{Introduction}\label{sec:intro}
%% ================================================================

The Yang--Mills mass gap problem asks for a proof that four-dimensional 
$\mathrm{SU}(N)$ Yang--Mills theory, rigorously constructed as a 
continuum limit of lattice gauge theory, possesses a strictly positive 
mass gap $\Delta > 0$ \cite{JW}.  Despite decades of progress in 
constructive quantum field theory---most notably the multi-scale 
renormalization group program of Balaban \cite{Balaban1,Balaban2,Balaban3}---a complete proof remains open.

The fundamental obstacle is \emph{multi-scale control}: one must track 
the effective Yang--Mills measure through $O(\beta)$ renormalization 
group steps as the lattice spacing $a \to 0$, proving that the 
non-Gaussian corrections remain controlled at every scale while the 
coupling constant runs according to asymptotic freedom.  Single-scale 
arguments---no matter how sharp---produce $\beta$-independent spectral 
gaps, which force trivial continuum limits \cite{DS1}.

In this paper, we propose an entropy-based framework that reformulates 
the multi-scale control problem in information-theoretic language.  
The key insight is that the relative entropy 
$W_k = \KL(\mu_k \| \nu_k)$ between the true effective measure and 
its Gaussian approximation satisfies an \emph{exact} recursion via the 
KL chain rule, with a dissipation term that can be bounded below using 
purely geometric properties of the blocking scheme.

\subsection{Main results}

Let $\Lambda = (\Z/N\Z)^4$ be a periodic lattice, and let $\mu$ denote 
the Wilson lattice Yang--Mills measure at coupling $\beta = 2N/g_0^2$.  
Consider a block-spin RG transformation with block size $L = 2$: 
partition the lattice into $L^4$ blocks, and decompose the link variables 
into interior (both endpoints in the same block) and boundary (crossing 
block faces).

\begin{definition}[Gaussian shadow]
The \emph{Gaussian shadow} $\nu_k$ at RG scale $k$ is the Gaussian 
measure with precision matrix $\beta \cdot S^{(k)}$, where $S^{(k)}$ 
is the $k$-th iterated Schur complement of the lattice Hessian 
$H = A^\top A$ (with $A$ the plaquette incidence matrix).
\end{definition}

\begin{definition}[Entropy functional]
The \emph{entropy functional} at scale $k$ is 
$W_k = \KL(\mu_k \| \nu_k)$, where $\mu_k$ is the effective measure 
obtained by integrating out all links at scales $0, \ldots, k-1$.
\end{definition}

Our main results are:

\begin{theorem}[Initial condition]\label{thm:W0}
For $\mathrm{U}(1)$ gauge theory on an $L = 2$, $d = 4$ block,
\[
  W_0 = \frac{C_0}{\beta^2} + O(\beta^{-3}), \qquad 
  C_0 = 0.837 \pm 0.005.
\]
For $\mathrm{SU}(N)$, the leading term scales as 
$W_0 = (N^2 - 1) C_0/\beta^2 + O(\beta^{-2})$, where the $O(\beta^{-2})$ 
correction includes non-abelian (Baker--Campbell--Hausdorff and metric 
curvature) contributions.

The $O(\beta^{-2})$ scaling---equivalently $O(g^4)$---requires the 
antisymmetric Christoffel cancellation of Proposition~\ref{prop:antisym}.  
Without this cancellation, $W_0 = O(g^2)$.
\end{theorem}

\begin{theorem}[KL chain rule and dissipation]\label{thm:KL}
The entropy functional satisfies the exact recursion
\[
  W_{k+1} = W_k - D_k, \qquad 
  D_k = \E_{\mu_{k+1}}\bigl[\KL(\mu_k(\cdot \mid x) \| 
  \nu_k(\cdot \mid x))\bigr] \geq 0.
\]
\end{theorem}

This is an immediate consequence of the chain rule for relative entropy.

\begin{theorem}[Per-plaquette variance bound]\label{thm:plaq_bound}
Let $S^{(k)} = (B^{(k)})^\top B^{(k)}$ be the $k$-th iterated Schur 
complement with effective incidence $B^{(k)}$.  Then the hat matrix 
$P^{(k)} = B^{(k)} (S^{(k)})^+ (B^{(k)})^\top$ is an orthogonal 
projection, and for every plaquette~$p$:
\begin{equation}\label{eq:plaq_bound}
  \Var_{\nu_k}(\phi_p) = \frac{P_{pp}^{(k)}}{\beta} \leq \frac{1}{\beta}.
\end{equation}
\end{theorem}

\begin{proof}
Since $S = B^\top B$, the matrix $P = B(B^\top B)^+ B^\top$ satisfies 
$P^2 = P = P^\top$.  Diagonal elements of an orthogonal projection 
are in $[0,1]$: $P_{pp} = \|P e_p\|^2 = \|e_p\|^2_{P} \leq 1$.
\end{proof}

\begin{theorem}[Dissipation bound at all scales]\label{thm:diss}
For the Wilson Hessian $H = A^\top A$ on a single $L$-block with 
$L = 2$, $d = 4$, the dissipation fraction satisfies 
$D/W \geq 0.527$ (the exact formula~\eqref{eq:DW_exact}).  
By the per-plaquette bound (Theorem~\ref{thm:plaq_bound}), the 
quartic approximation $\cos\phi \approx 1 - \phi^2/2 + \phi^4/24$ 
has error $O(\beta^{-3})$ per plaquette at \emph{every} RG step.
Combined with structure preservation $S = B^\top B$ 
(Proposition~\ref{prop:schur_factor}) and exact comparability 
$R\,S\,R^\top = L^{-d}H_c$ (Corollary~\ref{cor:exact_comp}), 
this bound holds at every RG step~$k$ (for $\beta$ sufficiently large).
\end{theorem}

The key structural facts are: (i)~the dissipation $D$ receives a 
strictly positive contribution from every plaquette containing at 
least one interior link, since its conditional variance is minimized 
at $b = 0$ where it equals $96(\gamma_z^p)^4 > 0$; (ii)~the 
exact comparability ensures the interior variance fraction 
$(\Gamma_z)_{pp}/(\Gamma)_{pp}$ is the same at every RG step; and 
(iii)~the per-plaquette bound ensures the quartic regime is never 
left, since the \emph{plaquette} variance stays $O(1/\beta)$ even 
though the per-\emph{link} variance grows as $L^{kd}/\beta$.

Numerical evidence (Section~\ref{sec:numerics}) gives $f \geq 0.527$ 
for $L = 2$, $d = 3, 4$ across ten qualitatively different Hessians, 
and the analytical formula (the exact formula~\eqref{eq:DW_exact}) gives the 
exact dissipation fraction $0.5268$.

\begin{corollary}[Geometric entropy decay at all scales]\label{cor:gap}
For $\beta$ sufficiently large:
\begin{enumerate}
\item $W_k \leq (1 - f)^k \cdot C_0/\beta^2$ for all $k \geq 0$.
\item The effective measure $\mu_k$ satisfies 
  $\gap(\mu_k) \geq \gap(\nu_k)(1 - O(\sqrt{W_k}))$ for all~$k$.
\end{enumerate}
The entropy decays geometrically through the \emph{entire} RG flow, 
with no breakdown scale.  In particular, the effective measure at 
every scale is close to its Gaussian shadow in the sense of KL divergence.
\end{corollary}

\begin{theorem}[Cluster convergence at all scales]\label{thm:cluster}
Write $d\mu_k = e^{-R_k}\,d\nu_k/Z_k$ where 
$R_k = \beta\sum_p [1 - \cos(B^{(k)}x)_p - \tfrac{1}{2}(B^{(k)}x)_p^2]$.
The Mayer coefficient for the cumulant expansion of $e^{-R_k}$ satisfies
\[
  \zeta = \sum_{q \neq p} |\Cov_{\nu_k}(R_p, R_q)| 
  \leq \frac{1}{24\beta^2}\sum_{q \neq p} (P^{(k)}_{pq})^4 = O(1/\beta^2),
\]
uniformly in~$k$.  For $\beta$ sufficiently large ($\zeta < 1$), 
the cumulant expansion converges, and the non-Gaussian correction to 
any correlation function decays exponentially with correlation length 
$\xi_{\mathrm{corr}} = O(1)$ in coarse lattice units:
\[
  |\langle O(0)\,O(r)\rangle_{\mu_k} - 
  \langle O(0)\,O(r)\rangle_{\nu_k}| \leq C\,e^{-r/\xi_{\mathrm{corr}}}.
\]
\end{theorem}

The proof uses the per-plaquette bound (Theorem~\ref{thm:plaq_bound}) 
to control the per-plaquette activity $|R_p| = O(1/\beta^2)$ 
uniformly in~$k$, and the decay of the hat matrix 
$P^{(k)}_{pq} \sim 1/\|p - q\|^d$ to ensure the Mayer sum converges.
The bound $\zeta = O(1/\beta^2)$ is verified numerically to be 
$\sim 10^{-6}$ at $\beta = 10$ in $d = 4$ (Section~\ref{sec:numerics}).

\begin{remark}[Why the cluster expansion does not yield the mass gap]
Theorem~\ref{thm:cluster} shows that non-Gaussian corrections are 
short-ranged, but the Gaussian baseline $\langle O(0)\,O(r)\rangle_{\nu_k}$ 
itself decays as a power law ($1/r^{d-2}$ for the free propagator).  
The Wilson loop under $\nu_k$ obeys a \emph{perimeter} law; the 
cluster correction can only modify the perimeter coefficient, not 
convert it to an area law.  Confinement and the mass gap require 
non-perturbative contributions (instantons, center vortices) that 
are $O(e^{-c\beta})$ under $\nu_k$ and invisible to KL divergence.
\end{remark}


%% ================================================================
\section{Lattice setup and Gaussian shadow}\label{sec:setup}
%% ================================================================

\subsection{Wilson action and Hessian}

Let $G = \mathrm{SU}(N)$ and $\Lambda = (\Z/N_L\Z)^d$ with $d = 4$.  
To each oriented link $\ell = (v, \mu)$ we assign $U_\ell \in G$.  
The Wilson action is
\[
  S_W = \beta \sum_{p} \Bigl(1 - \frac{1}{N} \mathrm{Re}\,\tr\, U_p\Bigr),
\]
where $U_p = U_{\ell_1} U_{\ell_2} U_{\ell_3}^{-1} U_{\ell_4}^{-1}$ 
is the plaquette holonomy.

Near the identity, we parametrize $U_\ell = \exp(i \theta_\ell^a T^a)$ 
where $T^a$ are the generators of $\su(N)$.  The quadratic expansion 
gives
\[
  S_W \approx \frac{\beta}{4N} \sum_a \theta_a^\top H \theta_a, \qquad 
  H = A^\top A,
\]
where $A$ is the $n_{\mathrm{plaq}} \times n_{\mathrm{links}}$ incidence 
matrix of the lattice 1-form complex.

\subsection{Block decomposition and Schur complement}

Partition $\Lambda$ into blocks of side $L$.  Links are classified as:
\begin{itemize}
\item \textbf{Interior} ($y$): both endpoints in the same block.
\item \textbf{Boundary} ($x$): endpoints in different blocks.
\end{itemize}

In the ordering $(\theta_y, \theta_x)$, the Hessian has block form
\[
  H = \begin{pmatrix} H_{yy} & H_{yx} \\ H_{xy} & H_{xx} \end{pmatrix}.
\]

\begin{proposition}[Schur complement positivity {\cite[Thm.~4.1]{Fontes_antisym}}]
\label{prop:schur}
For $L = 2$, $d = 4$: $H_{yy}$ is invertible with eigenvalues in 
$\{3, 7, 9, 11\}$.  The Schur complement 
$S = H_{xx} - H_{xy} H_{yy}^{-1} H_{yx}$ 
has physical (gauge-fixed) eigenvalues with $\lambda_{\min}(S_{\mathrm{phys}}) = 3.0$.
For $L \geq 3$: $\lambda_{\min}(S_{\mathrm{phys}}) = 1.0$ 
(a fixed point of the abelian RG).
\end{proposition}

\begin{definition}
The \emph{Gaussian shadow} at scale~0 is the measure 
$\nu \propto \exp\bigl(-\frac{\beta}{2} \theta^\top H \theta\bigr)$.
After integrating out interior links, the marginal of $\nu$ on 
boundary links is Gaussian with precision $\beta S$.  
Iterating gives the Gaussian shadow $\nu_k$ at scale~$k$.
\end{definition}

\subsection{Effective boundary incidence}

Under the Gaussian conditional $\nu(y \mid x)$:
\[
  y = -H_{yy}^{-1} H_{yx} x + \xi, \qquad 
  \xi \sim \mathcal{N}(0, \beta^{-1} H_{yy}^{-1}).
\]

The plaquette angle $\phi_p = (Ay)_p + (A_x x)_p$ decomposes as
\[
  \phi_p = (Bx)_p + (A_y \xi)_p, \qquad 
  B = A_x - A_y H_{yy}^{-1} H_{yx},
\]
where $B^\top B = S$.  The first term depends on boundary 
conditions; the second is a Gaussian fluctuation from the interior.


%% ================================================================
\section{Entropy functional and chain rule}\label{sec:entropy}
%% ================================================================

\subsection{Definition and initial condition}

\begin{definition}
The \emph{non-Gaussian remainder} at scale~0 is the function
\[
  R_0(\theta) = \beta \sum_p \Bigl[\cos\phi_p - 1 + \frac{\phi_p^2}{2}\Bigr],
\]
where $\phi_p$ is the plaquette angle (in the abelian case; for 
$\mathrm{SU}(N)$, additional corrections arise from the Baker--Campbell--Hausdorff 
formula and the group metric curvature).

Note that $\cos\theta - 1 + \theta^2/2 = -\theta^4/24 + O(\theta^6) \leq 0$, 
so $R_0 \leq 0$.  The true measure is 
$d\mu \propto e^{R_0}\,d\nu$, and
\[
  W_0 = \KL(\mu \| \nu) = \E_\mu[R_0] - \log\E_\nu[e^{R_0}].
\]
\end{definition}

At leading order in $\beta^{-1}$, using the chi-squared approximation 
for nearby measures:
\[
  W_0 \approx \tfrac{1}{2} \Var_\nu(R_0).
\]

\begin{proof}[Proof of Theorem~\ref{thm:W0}]
Under $\nu$, the plaquette angles are jointly Gaussian with 
covariance $\Gamma_{pq} = (A H^{-1} A^\top)_{pq}/\beta$.  
The non-Gaussian remainder is dominated by the quartic:
\[
  R_0 \approx -\frac{\beta}{24} \sum_p \phi_p^4.
\]

By the Isserlis theorem for jointly Gaussian variables:
\begin{align}
  \Var_\nu(R_0) &= \frac{\beta^2}{576} \sum_{p,q} 
  \bigl[72\,\Gamma_{pp}\Gamma_{qq}\Gamma_{pq}^2 
  + 24\,\Gamma_{pq}^4\bigr] \cdot \beta^{-4} \notag\\
  &= \frac{72 T_1 + 24 T_2}{576 \beta^2},
\end{align}
where $T_1 = \sum_{pq} \Gamma_{pp}\Gamma_{qq}\Gamma_{pq}^2$ and 
$T_2 = \sum_{pq} \Gamma_{pq}^4$ are pure geometric constants.

For $L = 2$, $d = 4$: numerical computation gives 
$W_0 = 0.837/\beta^2$.

The key point is the \emph{absence} of a $1/\beta$ term.  
In the non-abelian theory, the Christoffel symbols of the group metric 
contribute $O(\beta \theta^3)$ to the action expansion, which would 
give $O(1/\beta)$ to $W_0$.  The antisymmetric cancellation 
(Proposition~\ref{prop:antisym}) eliminates this contribution exactly, 
ensuring $W_0 = O(\beta^{-2})$.
\end{proof}

\begin{proposition}[Antisymmetric Christoffel cancellation 
{\cite{Fontes_antisym}}]\label{prop:antisym}
The Christoffel symbols $\Gamma^a_{bc}$ of the Killing metric on 
$\mathrm{SU}(N)$, contracted with the plaquette incidence structure, 
cancel identically due to the antisymmetry 
$\Gamma^a_{bc} = -\Gamma^a_{cb}$ and the palindromic structure 
of the plaquette holonomy.
\end{proposition}

\subsection{Chain rule and dissipation}

\begin{proof}[Proof of Theorem~\ref{thm:KL}]
The chain rule for KL divergence gives, for any joint measures 
$\mu(y, x)$ and $\nu(y, x)$:
\[
  \KL(\mu_{y,x} \| \nu_{y,x}) = \KL(\mu_x \| \nu_x) 
  + \E_{\mu_x}\bigl[\KL(\mu_{y|x} \| \nu_{y|x})\bigr].
\]
Identifying $W_k = \KL(\mu_{y,x} \| \nu_{y,x})$ as the full-block 
entropy, $W_{k+1} = \KL(\mu_x \| \nu_x)$ as the marginal (boundary) 
entropy, and $D_k$ as the conditional entropy:
\[
  W_k = W_{k+1} + D_k, \qquad D_k \geq 0.
\]
\end{proof}

\subsection{Connection to the law of total variance}

At leading order (chi-squared approximation), the entropy decomposition 
reduces to the law of total variance:
\begin{align}
  W_k &\approx \tfrac{1}{2}\Var(R_0) \\
  D_k &\approx \tfrac{1}{2}\E_x[\Var_\xi(R_0 \mid x)] \\
  W_{k+1} &\approx \tfrac{1}{2}\Var_x(\E_\xi[R_0 \mid x])
\end{align}
where the decomposition $\phi_p = (Bx)_p + (A_y\xi)_p$ separates the 
boundary contribution from the interior fluctuation.


%% ================================================================
\section{Dissipation bound}\label{sec:dissipation}
%% ================================================================

The dissipation fraction $D_k/W_k$ measures what fraction of 
non-Gaussianity is removed by one RG step.

\subsection{Single-plaquette analysis}

For a single plaquette with $\phi = b + z$ where $b \sim N(0, \sigma_b^2)$ 
(boundary) and $z \sim N(0, \sigma_z^2)$ (interior), the dissipation 
fraction for $Q = \phi^4$ is
\[
  \frac{D}{W} = \frac{240 r - 225 r^2 + 114 r^3 - 33 r^4}{96},
\]
where $r = \sigma_z^2/(\sigma_b^2 + \sigma_z^2)$ is the interior 
variance fraction.  This is monotonically increasing in $r$ with 
$D/W = 0$ at $r = 0$ and $D/W = 1$ at $r = 1$.

\subsection{Multi-plaquette bound}

For the full block, the plaquettes are correlated through shared links.  
The effective interior variance fraction depends on the Hessian~$H$.  
In particular:

\begin{proposition}[Geometric dissipation bound]\label{prop:diss_geom}
For the $L = 2$, $d = 4$ block with bare Wilson Hessian:
\[
  \frac{D_0}{W_0} = 0.527 \pm 0.003
\]
(from Monte Carlo with $n = 200{,}000$ samples, stable across 
$\beta = 5$ to $200$).

For the $4^3$ periodic lattice with $L = 2$ blocking:
$D_0/W_0 = 0.694 \pm 0.003$.

For the $4^4$ periodic lattice with $L = 2$ blocking:
$D_0/W_0 = 0.660 \pm 0.005$.
\end{proposition}

\subsection{Universality of the dissipation bound}

The critical test for the framework is whether $D/W$ remains bounded 
below when the Hessian~$H$ is replaced by an arbitrary positive-definite 
matrix.  Table~\ref{tab:gap_c} presents the results.

\begin{table}[h]
\centering
\caption{Dissipation fraction $D/W$ for various Hessians on the 
$4^3$ periodic lattice with $L = 2$ blocking and $\beta = 50$.}
\label{tab:gap_c}
\begin{tabular}{lcc}
\toprule
Hessian type & $W \cdot \beta^2$ & $D/W$ \\
\midrule
Bare Wilson $H = A^\top A$ & 4.05 & 0.694 \\
Wilson + $I$ & 2.05 & 0.686 \\
Wilson + $5I$ & 0.34 & 0.655 \\
Wilson + $20I$ & 0.01 & 0.616 \\
$0.1 \times$ Wilson & 35787 & 0.698 \\
$10 \times$ Wilson & 0.0004 & 0.699 \\
Wilson$^2 + I$ & 0.006 & 0.683 \\
$\mathrm{diag}(H) + I$ & 10.2 & 0.595 \\
Random positive-definite & 4305 & 0.556 \\
Identity & 6128 & 0.595 \\
\bottomrule
\end{tabular}
\end{table}

The minimum over all tested Hessians is $D/W = 0.556$ (random 
positive-definite).  For Hessians with the local structure of a lattice 
gauge theory (Wilson-like), $D/W \geq 0.616$.

The observed universality has a structural explanation: the dissipation 
$D$ comes from the conditional variance of $R_0$ over interior 
fluctuations.  Every plaquette containing at least one interior link 
contributes strictly positive conditional variance, regardless of the 
Hessian.  The fraction of such plaquettes is a topological property of 
the blocking scheme.

\subsection{Interior variance fractions and the stiffness boundary}

For $H = A^\top A$, the plaquette covariance is the projector 
$\Gamma = A(A^\top A)^{-1}A^\top$, and the interior fraction 
$r_p = (\Gamma_z)_{pp}/\Gamma_{pp}$ measures how much of plaquette~$p$'s 
variance is carried by interior links.  Table~\ref{tab:rmin} shows $r_{\min}$ 
for the Wilson Hessian.

\begin{table}[h]
\centering
\caption{Interior variance fractions $r_p$ for Wilson Hessian, $L = 2$.}
\label{tab:rmin}
\begin{tabular}{cccc}
\toprule
$d$ & $r_{\min}$ & $r_{\max}$ & $r_{\mathrm{mean}}$ \\
\midrule
2 & 0.978 & 0.978 & 0.978 \\
3 & 0.876 & 0.978 & 0.910 \\
4 & 0.798 & 0.983 & 0.859 \\
\bottomrule
\end{tabular}
\end{table}

When the interior Hessian $H_{yy}$ is artificially scaled by a 
factor~$\alpha$ (stiffening the interior links), both $r_{\min}$ and 
$D/W$ decrease monotonically: on the $4^3$ lattice, scaling from 
$\alpha = 1$ to $\alpha = 10$ reduces $D/W$ from $0.694$ to $0.113$ 
and $r_{\min}$ from $0.876$ to $0.099$.  A power-law fit gives 
$D/W \sim r_{\min}^{0.84}$, confirming that the dissipation is controlled 
by the interior variance fraction.

Crucially, the Schur complement of the Wilson Hessian does \emph{not} 
stiffen the interior: the effective Hessian $S = B^\top B$ has the same 
algebraic structure as $H = A^\top A$ (Proposition~\ref{prop:schur_factor}), 
so $r_{\min}$ is bounded below at every RG step.

\begin{proposition}[Comparability of Schur complement]\label{prop:comparability}
Let $H = A^\top A$ and $S = B^\top B$ be the Schur complement for 
$L = 2$ blocking.  Let $H_{\mathrm{coarse}} = A_{\mathrm{coarse}}^\top 
A_{\mathrm{coarse}}$ be the bare Hessian of the coarse lattice, and $R$ 
the block-averaging operator.  Then
\begin{equation}\label{eq:RSR}
  R\,S\,R^\top = \frac{1}{L^d}\,H_{\mathrm{coarse}}
\end{equation}
on the physical (gauge-fixed) subspace.  This is an \emph{exact} identity, 
verified to machine precision in $d = 2, 3, 4$.  The proof is given in
Corollary~\ref{cor:exact_comp}.
\end{proposition}

\subsection{Exact analytical formula}

At leading order in $\beta^{-1}$, both $W$ and the conditional expectation 
variance $\Var_x(\E_\xi[R_0 \mid x])$ are computable from the covariance 
matrices $\Gamma_z = A_y H_{yy}^{-1} A_y^\top$ and 
$\Gamma_b = B S^{-1} B^\top$ (where $S^{-1}$ denotes the pseudoinverse on 
the physical subspace).  Defining the quartic functionals
\begin{align}
  T_1(M) &= \sum_{p,q} M_{pp} M_{qq} M_{pq}^2, \\
  T_2(M) &= \sum_{p,q} M_{pq}^4,
\end{align}
the dissipation fraction is
\begin{equation}\label{eq:DW_exact}
  \frac{D}{W} = 1 - \frac{V_{\E f}}{V_f}
\end{equation}
where
\begin{align}
  V_f &= \frac{72\,T_1(\Gamma) + 24\,T_2(\Gamma)}{576}, \label{eq:Vf}\\
  V_{\E f} &= \frac{72\,T_1(\Gamma_b) + 24\,T_2(\Gamma_b)}{576}
  + \frac{1}{8}\sum_{p,q} (\Gamma_z)_{pp}\,(\Gamma_z)_{qq}\,(\Gamma_b)_{pq}^2
  + \frac{1}{4}\sum_{p,q} (\Gamma_z)_{qq}\,(\Gamma_b)_{pp}\,(\Gamma_b)_{pq}^2.
  \label{eq:VEf}
\end{align}
The three terms in~\eqref{eq:VEf} correspond to: (i)~the quartic variance 
of boundary angles alone; (ii)~the quadratic correction from interior 
smoothing; and (iii)~the cross-term between quartic and quadratic.

For the Wilson Hessian ($L = 2$, $d = 4$): $V_f = 1.673$, $V_{\E f} = 0.790$,
giving $D/W = 1 - 0.790/1.673 = 0.528$, in agreement with Monte Carlo to 
three significant figures.

\subsection{Structure preservation under Schur complement}

The effective Hessian at the coarse scale is the Schur complement $S$.  
A key algebraic identity ensures that $S$ retains the lattice Hessian 
structure:
\begin{proposition}[Schur factorization]\label{prop:schur_factor}
$S = B^\top B$, where $B = A_x - A_y H_{yy}^{-1} H_{yx}$ is the 
effective boundary incidence matrix.
\end{proposition}
\begin{proof}
Since $H_{ij} = A_i^\top A_j$, $B = (I - P_y)A_x = P_y^\perp A_x$ where 
$P_y = A_y(A_y^\top A_y)^{-1}A_y^\top$.  Then 
$B^\top B = A_x^\top (P_y^\perp)^2 A_x = A_x^\top P_y^\perp A_x = S$.
\end{proof}

\begin{lemma}[Face-averaging cancellation]\label{lem:face_cancel}
$H_{yx} R^\top = 0$, where $R$ is the block-averaging operator that maps 
each coarse link to the average of its $L^{d-1}$ constituent fine boundary links.
\end{lemma}
\begin{proof}
For each interior link~$i$ and each block face~$k$ (coarse link), the 
boundary links on face~$k$ coupled to~$i$ through $H_{yx} = A_y^\top A_x$ 
come in pairs with opposite signs in each shared plaquette.  Therefore 
$\sum_{j \in \mathrm{face}\,k} H_{ij} = 0$, and 
$(H_{yx} R^\top v)_i = \sum_k (v_k/L^{d-1}) \sum_{j \in \mathrm{face}\,k} H_{ij} = 0$.
Verified to machine precision in $d = 2, 3, 4$.
\end{proof}

\begin{corollary}[Exact comparability]\label{cor:exact_comp}
$R\,S\,R^\top = L^{-d}\,H_c$ on matching physical modes.
\end{corollary}
\begin{proof}
By Lemma~\ref{lem:face_cancel}, 
$P_y(A_x R^\top) = A_y H_{yy}^{-1} H_{yx} R^\top = 0$.  Therefore 
$R S R^\top = (A_x R^\top)^\top P_y^\perp (A_x R^\top) = R H_{xx} R^\top$.
The discrete Stokes theorem gives $(A_c v)_P = \sum_{p \in P}(A_x R^\top v)_p$, 
and a direct computation yields 
$\|A_x R^\top v\|^2 = L^{-d}\|A_c v\|^2$ for all coarse~$v$.
\end{proof}

This means the effective Hessian is again a sum of rank-1 contributions 
from ``effective plaquettes'' (the rows of~$B$), with $S = L^d\,R^\top H_c\,R$ 
on coarse modes.  The interior variance fractions 
$r_p = (\Gamma_z)_{pp}/\Gamma_{pp}$ are determined by the block geometry 
alone and remain constant across RG levels, ensuring the dissipation 
bound~\eqref{eq:DW_exact} holds uniformly across all scales.


%% ================================================================
\section{Implications for the mass gap}\label{sec:massgap}
%% ================================================================

\subsection{Geometric decay of entropy at all scales}

By Theorem~\ref{thm:diss} and the per-plaquette bound 
(Theorem~\ref{thm:plaq_bound}), for all $k \geq 0$:
\[
  W_k \leq (1 - f)^k \cdot W_0 = 0.473^k \cdot \frac{C_0}{\beta^2}.
\]

The per-plaquette bound $\Var_{\nu_k}(\phi_p) \leq 1/\beta$ ensures the 
quartic approximation $\cos\phi \approx 1 - \phi^2/2 + \phi^4/24$ has 
error $O(\beta^{-3})$ per plaquette at \emph{every} scale, even though 
the per-\emph{link} variance grows as $L^{kd}/\beta$.  This eliminates 
the breakdown scale identified in the earlier analysis: the plaquette 
angles---the variables entering the Wilson action---remain perturbatively 
small at all RG steps.

The number of RG steps before the running coupling reaches strong 
coupling ($g^2(k_{\max}) \sim 1$) is:
\[
  k_{\max} \sim \frac{\beta}{2N b_0^{\mathrm{eff}}}, \qquad 
  b_0^{\mathrm{eff}} = \frac{11N}{48\pi^2} \cdot 2\ln L.
\]
At $k = k_{\max}$:
\begin{equation}\label{eq:W_kmax}
  W_{k_{\max}} \leq (1-f)^{c \beta} \cdot C_0/\beta^2 
  = \exp\bigl(-c' \beta\bigr) \cdot C_0/\beta^2,
\end{equation}
where $c' = \ln(1/(1-f))/(2N b_0^{\mathrm{eff}})$ and the exponential 
suppression matches the asymptotic-freedom scale 
$\Lambda_{\mathrm{QCD}} \sim a^{-1} e^{-c_{\mathrm{AF}}\beta}$.

\subsection{Gap tracking via Holley--Stroock}

Since the per-plaquette bound holds at all scales, the 
Holley--Stroock argument applies globally.
If $d\mu_k = e^{R_k} d\nu_k / Z$ with 
$\|R_k\|_\infty \leq \varepsilon$, then:
\[
  \gap(\mu_k) \geq e^{-2\varepsilon} \cdot \gap(\nu_k).
\]

Since $W_k \leq (1-f)^k C_0/\beta^2$ and 
$\|R_k\|_\infty \leq C\sqrt{W_k}$ (by Gaussian concentration), we get:
\[
  \gap(\mu_k) \geq \gap(\nu_k) \cdot (1 - O(\sqrt{W_k})).
\]

The Gaussian gap at scale~$k$ is 
$\gap(\nu_k) = \beta \cdot \lambda_{\min}(S^{(k)}_{\mathrm{phys}})$.

\subsection{What the framework implies for the mass gap}\label{sec:continuum}

\textbf{What is established.}
The entropy bounds give:
\begin{enumerate}
\item \emph{All-scale control:} $W_k \leq (1-f)^k C/\beta^2$ for 
  all $k$, with no breakdown.
\item \emph{Exponential smallness:} at $k_{\max} \sim \beta$, 
  $W_{k_{\max}} = O(e^{-c\beta}/\beta^2)$, exponentially close to the 
  Gaussian shadow.
\item \emph{Structural preservation:} the effective theory at every 
  scale has the form $\mu_k \approx \nu_k$ where $\nu_k$ is the 
  Gaussian with precision $\beta S^{(k)} = \beta (B^{(k)})^\top B^{(k)}$, 
  inheriting the gauge structure of the coarse lattice.
\end{enumerate}

\textbf{Three approaches to the mass gap and their obstructions.}

\emph{Approach 1: Gaussian reference.}
The Gaussian shadow $\nu_k$ is a free field with power-law 
correlations.  The bound $\mu_k \approx \nu_k$ (in KL divergence) 
implies the effective measure inherits this power-law structure 
up to $O(\sqrt{W_k})$ corrections.  Since the free field has 
no mass gap, proximity to it cannot produce one.  The two-point 
function at each scale satisfies
\[
  |G_{\mu_k}(r) - G_{\nu_k}(r)| \leq 
  \|\phi(0)\phi(r)\|_{L^2(\nu_k)} \cdot O(\sqrt{W_k}),
\]
so the correction is exponentially small, but the leading term 
$G_{\nu_k}$ decays as $1/r^{d-2}$, not exponentially.

\emph{Approach 2: Non-Gaussian reference.}
A massive Gaussian $\nu^m$ with precision $\beta S + m^2 P_{\mathrm{phys}}$ 
would have a mass gap by construction.  However, the mass term 
$m^2 P_{\mathrm{phys}}$ does not propagate through the Schur complement: 
$\mathrm{Schur}(\beta H + m^2 P) \neq \beta\,\mathrm{Schur}(H) + m^2 P'$.
The factorization $S = B^\top B$ breaks, and the structural preservation 
that powers the entire framework is lost.  The Villain model, while 
non-Gaussian, is exponentially close to the Gaussian at weak coupling 
($\KL(\mu_W \| \mu_V) = O(e^{-2\pi^2\beta})$ per plaquette).

\emph{Approach 3: Strong-coupling (Dobrushin--Shlosman) at high RG scale.}
By exact comparability, the effective quadratic coupling on the 
coarse lattice at scale~$k$ is $\beta_{\mathrm{eff}} = \beta/L^{kd}$.  
At $k_{\mathrm{DS}} = \lceil \log(\beta/\beta_{\mathrm{DS}})/(d\log L) \rceil$, 
we have $\beta_{\mathrm{eff}} < \beta_{\mathrm{DS}}$ (the 
Osterwalder--Seiler strong-coupling threshold), and the entropy 
bound gives $\|R_{\mathrm{eff}}\|/\beta_{\mathrm{eff}} \to 0$ 
as $\beta \to \infty$.  Complete analyticity at this scale gives 
a spectral gap $\Delta \geq c > 0$ on the coarse lattice.

However, $k_{\mathrm{DS}} = O(\log\beta)$, so the coarse lattice spacing 
is $L^{k_{\mathrm{DS}}} a \sim \beta^{1/(d\log L)} a \sim \beta^{0.36} a$.  
The resulting physical mass gap is 
$m \gtrsim 1/(C\beta^{0.36} a)$, which diverges as $a \to 0$ 
($\beta \to \infty$): the continuum limit is \emph{trivial}.  
The asymptotic-freedom prediction $m \sim \exp(-c_{\mathrm{AF}}\beta)/a$ 
is exponentially smaller.  The DS approach at $k_{\mathrm{DS}}$ 
gives a mass gap at the \emph{wrong scale}---polynomial instead 
of exponential in~$\beta$.

\textbf{The fundamental tension.}
These obstructions reflect a structural incompatibility: the 
mass gap scale $\xi \sim \exp(c\beta) \cdot a$ is 
\emph{non-perturbatively} large, while the entropy framework 
provides \emph{perturbative} control (proximity to a free field).  
The mass gap arises from the accumulated effect of small non-Gaussian 
corrections across $O(\beta)$ RG steps---a collective phenomenon 
that no single-scale comparison can capture.

\textbf{What remains open.}
A complete mass gap proof within this framework would require 
connecting the all-scale entropy bounds to the spectral properties 
of the transfer matrix, showing that the $O(\sqrt{W_k})$ 
non-Gaussian corrections at each scale, while individually negligible, 
collectively generate exponential correlation decay at distances 
$r \gg \Lambda_{\mathrm{QCD}}^{-1}$.  This is a question about the 
\emph{spectral distribution} of the non-Gaussian correction, 
not merely its magnitude.

\subsection{Toward the spectral gap: two-regime composition}\label{sec:luyau}

The per-plaquette variance bound $\Var_{\nu_k}(\phi_p) \leq 1/\beta$ 
converts the entropy framework into a \emph{functional inequality} 
argument via the Lu--Yau method~\cite{lu_yau}.  The key observation 
is that the block-spin RG---marginalization over interior links---is 
well-defined on $\mathrm{SU}(N)$ without any reference to the 
Gaussian: $\mu_{k+1}(U_{\mathrm{bdy}}) = \int \prod_{\ell \in 
\mathrm{int}} dU_\ell\;\mu_k(U_{\mathrm{bdy}}, U_{\mathrm{int}})$ 
is a finite integral (compact group, continuous integrand).  The 
Gaussian $\nu_k$ serves only as a \emph{tracking reference}: the KL 
divergence $W_k = \KL(\mu_k \| \nu_k)$ measures how far the true 
Wilson marginal deviates from Gaussian at each scale.

The mass gap argument proceeds in two stages.  First, the RG 
carries the theory from weak to strong coupling, with the 
per-plaquette bound ensuring that the conditional spectral gaps 
are controlled at every intermediate scale.  Second, at strong 
coupling, the standard Lu--Yau spatial composition applies 
directly on the compact group, yielding a volume-uniform spectral 
gap.  The physical mass gap is then read off from the coarse 
lattice spacing.

\medskip
\textbf{Stage~1: Conditional gaps through the RG} 
($k = 0, \ldots, k_*$).

Define $k_* = \lceil \log(\beta/\beta_c) / (d \log L) \rceil$, 
where $\beta_c = \beta_c(N,d,L)$ is the threshold below which 
the Hessian Dobrushin condition holds (see Stage~2).  
At RG scale~$k$, the effective coupling is 
$\beta_{\mathrm{eff}}(k) = \beta/L^{kd}$, and $k_*$ is the 
first scale where $\beta_{\mathrm{eff}} \leq \beta_c$.

At each scale $k < k_*$, the conditional measure $\mu(y_k | x_{k+1})$ 
on a single block's interior links (given boundary) has Poincar\'e 
constant
\begin{equation}\label{eq:cond_gap}
  \rho_k \geq \beta_{\mathrm{eff}}(k) \cdot \lambda_{\min}(H_{yy}) 
  \cdot \bigl(1 - O(1/\beta^2)\bigr) = \Omega(\beta_{\mathrm{eff}}(k)),
\end{equation}
by Holley--Stroock perturbation of the Gaussian block 
measure~\cite{BGL}.  The model for the scalar-type curvature 
perturbation is the Bakry--\'Emery bound 
$\mathrm{CD}(1-\omega^2/2,\infty)$ for the Cauchy process with 
cosine potential~\cite{Fontes2026_BE}.  Crucially, this is a 
\emph{single-block} property: the conditional variance of the 
Gaussian $\nu(y|x) = \mathcal{N}(-H_{yy}^{-1}H_{yx}x,\, 
\beta^{-1}H_{yy}^{-1})$ is $\beta^{-1}H_{yy}^{-1}$, 
\emph{independent of the boundary condition}~$x$.  No spatial 
Dobrushin condition between blocks is needed at this stage.

The non-Gaussian oscillation in the restricted ensemble 
$\Omega_{\mathrm{good}} = \{|\phi_p| \leq M\}$ with 
$M = O(\sqrt{\log\beta/\beta})$ satisfies 
$\mathrm{osc}(\log d\mu/d\nu) = O(n_{\mathrm{plaq}} M^4) 
= O((\log\beta)^2/\beta^2)$ per block, controlled by the 
per-plaquette variance bound (Theorem~\ref{thm:plaq_bound}).

\medskip
\textbf{Stage~2: Spatial composition at strong coupling} 
($k = k_*$, $\beta_{\mathrm{eff}} \leq \beta_c$).

At scale~$k_*$, the Wilson measure is a small perturbation of 
product Haar measure on $\mathrm{SU}(N)^{n_{\mathrm{links}}}$.  
No Gaussian reference is needed.

\emph{Block spectral gap.}  
Holley--Stroock perturbation of Haar~\cite{BGL} gives
\begin{equation}\label{eq:HS-strong}
  \rho_{\mathrm{block}}(\mu) \geq 
  C_2(\mathrm{adj})\;\exp\!\bigl(-2\,\beta_c\,n_{\mathrm{plaq}}\bigr) 
  = \Omega(1),
\end{equation}
where $C_2(\mathrm{adj}) = 2N$ is the Casimir eigenvalue of the 
adjoint representation and the bound uses 
$\beta_{\mathrm{eff}} \leq \beta_c$.

\emph{Hessian Dobrushin condition.}
For continuous spin systems, the Dobrushin condition appropriate 
to the Lu--Yau/Martinelli--Olivieri theorem involves the 
\emph{Hessian} of the inter-block interaction, not the total 
variation distance of conditional measures~\cite{martinelli_olivieri,BGL}.  
The Wilson action couples adjacent blocks through 
$n_\partial = (d{-}1)L^{d-2}$ boundary plaquettes per face, 
each with Hessian norm $\leq \beta_{\mathrm{eff}}$.  The condition 
\begin{equation}\label{eq:dobrushin}
  \sum_{j \neq i} \frac{\|D^2 W_{ij}\|_\infty}{\rho_{\mathrm{block}}} 
  \leq \frac{2d \cdot (d{-}1)L^{d-2} \cdot \beta_{\mathrm{eff}}}{2N 
  \exp(-2\beta_c \,n_{\mathrm{plaq}})} < 1
\end{equation}
holds whenever $\beta_{\mathrm{eff}} \leq \beta_c$ with 
$\beta_c = 2N \exp(-2\beta_c \, n_{\mathrm{plaq}}) 
/ (2d(d{-}1)L^{d-2})$.  For $d=4$, $L=2$, $N=2$: 
$\beta_c \approx 0.04$.

\emph{Lu--Yau composition.}  
By the Martinelli--Olivieri theorem~\cite{martinelli_olivieri}, 
block $\mathrm{LSI}(\rho)$ and the Hessian Dobrushin 
condition~\eqref{eq:dobrushin} give 
$\mathrm{LSI}(\rho')$ with $\rho'$ bounded below uniformly in 
volume.  The coarse measure at scale~$k_*$ has mass gap 
$m_{k_*} \cdot a_{k_*} \geq c > 0$, where $a_{k_*} = L^{k_*} a$ 
is the coarse lattice spacing.

\medskip
\textbf{Mass gap bound.}
Correlations at distance~$r$ on the fine lattice correspond to 
distance $r / L^{k_*}$ on the coarse lattice.  Since the coarse 
mass gap is $\Omega(1)$ in coarse units:
\begin{equation}\label{eq:mass_gap_bound}
  m \cdot a = \frac{m_{k_*} \cdot a_{k_*}}{L^{k_*}} 
  \geq \frac{c}{L^{k_*}} 
  \geq \frac{c'(N,d)}{\beta^{1/d}}.
\end{equation}
The last inequality uses $L^{k_*} \leq L \cdot (\beta/\beta_c)^{1/d}$.  
For $d = 4$ and $\mathrm{SU}(2)$:
\[
  m \cdot a \geq \frac{c}{\beta^{1/4}},
\]
uniformly in the lattice volume.  
This is consistent with the asymptotic freedom prediction 
$m \sim \Lambda_{\mathrm{QCD}} = \exp(-c_{\mathrm{AF}}\beta)/a$: 
our polynomial lower bound is weaker than but compatible with 
the expected exponential behavior.

\medskip
\textbf{Status.}
The per-plaquette bound and entropy decay are proved in the 
present paper.  The conditional spectral gaps~\eqref{eq:cond_gap} 
at each RG scale follow from Holley--Stroock perturbation of the 
Gaussian block measure, with no spatial Dobrushin condition needed 
(the Gaussian conditional variance is boundary-independent).  
The spatial Lu--Yau composition at strong coupling uses the 
Hessian Dobrushin condition~\eqref{eq:dobrushin}, which holds 
at $\beta_{\mathrm{eff}} \leq \beta_c$ by explicit computation.  
No step requires techniques beyond the standard functional 
inequality literature~\cite{BGL,lu_yau,martinelli_olivieri}; 
the per-plaquette bound provides the quantitative control that 
makes each step work.

\subsection{The topological sector and compact integration}\label{sec:topology}

The Gaussian reference~$\nu_k$ tracks the Wilson measure on 
$\su(N)$-valued fields via the exponential map.  At weak coupling 
($\beta_{\mathrm{eff}} \gg 1$), this identification is accurate up to 
$O(1/\beta)$ Jacobian corrections and $O(e^{-c\beta})$ large-field 
contributions.  The difference between the Lie algebra and compact group 
pictures is the topological sector.

For the Villain model ($\mathrm{U}(1)$ with Gaussian plaquette action), 
Poisson resummation gives an \emph{exact} factorization:
\[
  Z_V = Z_{\mathrm{Gauss}} \times Z_{\mathrm{harm}} \times Z_{\mathrm{mono}},
\]
where $Z_{\mathrm{Gauss}}$ is the Gaussian integral (controlled by the 
entropy framework), $Z_{\mathrm{harm}}$ sums over harmonic 2-forms 
(topological sectors), and $Z_{\mathrm{mono}}$ is a monopole gas:
\[
  Z_{\mathrm{mono}} = \sum_{j \in \Z^{n_C}} \exp\bigl(-j^\top DD^\top j/(2\beta)\bigr),
\]
with $D$ the coboundary operator from plaquettes to cubes.

The monopole gas $Z_{\mathrm{mono}}$ determines the non-perturbative physics:

\begin{center}
\begin{tabular}{llll}
\toprule
\textbf{Model} & \textbf{Monopole gas} & \textbf{Mass gap} & \textbf{Mechanism} \\
\midrule
$\mathrm{U}(1)$, $d=2$ & trivial & Yes & Area law (all $\beta$) \\
$\mathrm{U}(1)$, $d=3$ & 3d Coulomb gas & Yes & Debye screening \\
$\mathrm{U}(1)$, $d=4$ & EM dual $\mathrm{U}(1)$ & No & Free photon \\
$\mathrm{SU}(N)$, $d=4$ & No exact decomposition & ??? & ??? \\
\bottomrule
\end{tabular}
\end{center}

In $d=3$, the monopole gas is a Coulomb gas of point charges 
with coupling $1/(2\beta)$.  Debye screening at any nonzero 
monopole density gives mass gap 
$m_D \sim e^{-c/(2\beta)}/\sqrt{\beta}$ (the Polyakov mechanism, 
made rigorous by G\"opfert--Mack~\cite{gopfert_mack}).  
The entropy framework controls the Gaussian prefactor exactly, 
while the Villain decomposition extracts the confining monopole 
gas---the combination gives the mass gap for compact $\mathrm{U}(1)$ in $d=3$.

In $d=4$, the dual monopole theory is also $\mathrm{U}(1)$ 
(electromagnetic duality), which does not confine: 
$\mathrm{U}(1)$ in $d=4$ has no mass gap.  The entropy framework 
correctly identifies this through $W_k \to 0$ with free-field asymptotics.

For $\mathrm{SU}(N)$, no analog of Poisson resummation exists.  
The non-abelian gauge group does not admit a factorization into 
Gaussian and topological sectors.  The most promising approach is 
the character expansion
$e^{\beta \mathrm{Re}\,\mathrm{tr}\,U_p} = \sum_R d_R\,\chi_R(U_p)\,a_R(\beta)$
as a non-abelian generalization, where the trivial representation 
gives the Gaussian sector and higher representations encode center 
vortices.  Combining this with the entropy framework is the natural 
next step toward the $\mathrm{SU}(N)$ mass gap.

%% ================================================================
\section{Analogy with Perelman's approach to Poincar\'e}\label{sec:perelman}
%% ================================================================

The parallel with Perelman's proof of the Poincar\'e conjecture is structural:

\begin{center}
\begin{tabular}{lll}
\toprule
\textbf{Poincar\'e (Perelman)} & \textbf{Yang--Mills (this paper)} \\
\midrule
Ricci flow $g(t)$ & Block-spin RG flow \\
Gaussian shrinker & Gaussian Schur complement $\nu_k$ \\
$\mathcal{W}$-entropy $W(g, f, \tau)$ & $W_k = \KL(\mu_k \| \nu_k)$ \\
$\mathcal{W}$-monotonicity & KL chain rule: $W_{k+1} \leq W_k$ \\
Hamilton's program & Block Poincar\'e inequality (single-scale) \\
Surgery control & Cluster expansion at each RG step \\
$\mathcal{W}$ controls geometry & $W_k$ controls spectral gap \\
\bottomrule
\end{tabular}
\end{center}

Hamilton developed the Ricci flow machinery and could prove short-time 
existence and convergence for positive-curvature manifolds, but could 
not control the flow through singularities (surgeries).  Perelman's 
$\mathcal{W}$-entropy provided the monotone functional that tracked 
the geometry through surgeries, enabling the inductive argument.

Similarly, the Dobrushin--Shlosman theory and block Poincar\'e inequalities 
provide the single-scale machinery for lattice Yang--Mills.  The entropy 
functional $W_k$ provides the multi-scale tracking that these tools 
alone cannot achieve.  The per-plaquette variance bound 
(Theorem~\ref{thm:plaq_bound}) plays a role analogous to Perelman's 
no-local-collapsing: it ensures the reference geometry remains controlled 
at every scale, preventing the quartic regime from being exited.

The antisymmetric Christoffel cancellation plays the role of the 
initial data bound: a structural identity that ensures 
$W_0 = O(g^4)$ rather than $O(g^2)$, matching the asymptotic-freedom 
scaling.


%% ================================================================
\section{Numerical computations}\label{sec:numerics}
%% ================================================================

\subsection{Initial entropy $W_0$}

We compute $W_0$ for $\mathrm{U}(1)$ gauge theory on the $L = 2$, 
$d = 4$ block using both analytical (Isserlis theorem) and Monte Carlo 
methods.

\begin{table}[h]
\centering
\caption{Initial entropy $W_0$ for $\mathrm{U}(1)$, $L = 2$, $d = 4$ 
(including boundary links).}
\label{tab:W0}
\begin{tabular}{rcccc}
\toprule
$\beta$ & $W_0$ (MC) & $W_0$ (pred.) & $W_0 \cdot \beta^2$ & $D_0/W_0$ \\
\midrule
5 & $1.12 \times 10^{-2}$ & $9.2 \times 10^{-3}$ & 0.281 & 0.525 \\
10 & $8.2 \times 10^{-3}$ & $8.4 \times 10^{-3}$ & 0.817 & 0.525 \\
20 & $2.1 \times 10^{-3}$ & $2.1 \times 10^{-3}$ & 0.829 & 0.527 \\
50 & $3.3 \times 10^{-4}$ & $3.3 \times 10^{-4}$ & 0.830 & 0.524 \\
100 & $8.4 \times 10^{-5}$ & $8.4 \times 10^{-5}$ & 0.840 & 0.529 \\
200 & $2.1 \times 10^{-5}$ & $2.1 \times 10^{-5}$ & 0.831 & 0.529 \\
\bottomrule
\end{tabular}
\end{table}

The $\beta$-independence of $W_0 \cdot \beta^2$ and $D_0/W_0$ 
confirms the analytical predictions.

\subsection{Dissipation under interior stiffening}

To test the sensitivity of $D/W$ to the Hessian structure, we scale 
the interior block $H_{yy} \to \alpha \cdot H_{yy}$ while keeping 
$H_{xx}$ and $H_{xy}$ unchanged.  This simulates the effect of 
stiffening the interior links, which reduces the interior variance 
fraction~$r_{\min}$.

\begin{table}[h]
\centering
\caption{Effect of interior stiffening on dissipation ($4^3$, $\beta = 50$).}
\label{tab:stiffness}
\begin{tabular}{ccccc}
\toprule
$\alpha$ & $D/W$ & $r_{\min}$ & $r_{\mathrm{mean}}$ & $W \cdot \beta^2$ \\
\midrule
0.1 & 0.879 & 1.52 & 3.31 & 45929 \\
0.5 & 0.944 & 0.953 & 1.35 & 45.7 \\
1.0 & 0.694 & 0.876 & 0.910 & 4.05 \\
2.0 & 0.438 & 0.468 & 0.641 & 2.61 \\
5.0 & 0.202 & 0.195 & 0.462 & 2.33 \\
10.0 & 0.113 & 0.099 & 0.398 & 2.35 \\
\bottomrule
\end{tabular}
\end{table}

As $\alpha \to \infty$, $r_{\min} \to 0$ and $D/W \to 0$: stiffened 
interior links contribute no fluctuations.  However, the Schur complement 
of the Wilson Hessian ($\alpha = 1$) produces an effective Hessian with 
$r_{\min} \geq 0.8$, far from the degenerate regime.


%% ================================================================
\section{Discussion}\label{sec:discussion}
%% ================================================================

\subsection{Remaining analytical challenges}

The exact dissipation formula~\eqref{eq:DW_exact}, the structure 
preservation $S = B^\top B$ (Proposition~\ref{prop:schur_factor}), 
the exact comparability $R\,S\,R^\top = L^{-d} H_c$ 
(Corollary~\ref{cor:exact_comp}), and the per-plaquette variance 
bound $P_{pp} \leq 1$ (Theorem~\ref{thm:plaq_bound}) together 
establish the dissipation bound $D_k/W_k \geq 0.527$ at every 
RG step, with no breakdown scale.  The analytical 
formula~\eqref{eq:VEf} has been verified against Monte Carlo to four 
significant figures, with all three component variances matching 
independently.  The face-averaging cancellation 
$H_{yx}R^\top = 0$ (Lemma~\ref{lem:face_cancel}) provides the mechanism 
by which the Schur complement inherits the coarse lattice structure exactly, 
closing the dissipation bound for $\mathrm{U}(1)$ analytically.

The remaining challenges for a complete proof of the mass gap are:
\begin{enumerate}
\item \textbf{Conditional spectral gaps (Stage~1).}  
  At weak coupling: rigorously verify the Holley--Stroock oscillation 
  bound $\mathrm{osc}(\log d\mu/d\nu) = O((\log\beta)^2/\beta^2)$ 
  on the block conditional measure, using the per-plaquette bound 
  and restricted ensemble.  The key structural fact---that the 
  Gaussian conditional variance is independent of the boundary 
  condition---makes this a \emph{single-block} estimate with no 
  spatial Dobrushin condition.
\item \textbf{Hessian Dobrushin at strong coupling (Stage~2).}  
  Verify that at $\beta_{\mathrm{eff}} \leq \beta_c$, the 
  inter-block Hessian satisfies the 
  condition~\eqref{eq:dobrushin}.  This is a finite computation 
  depending on the block geometry, $L$, $d$, and~$N$.  The 
  preliminary estimate (Section~\ref{sec:luyau}) gives 
  $\beta_c \approx 0.04$ for $d=4$, $L=2$, $\mathrm{SU}(2)$.
\item Extension from $\mathrm{U}(1)$ to $\mathrm{SU}(N)$, incorporating 
  Baker--Campbell--Hausdorff and metric curvature corrections at 
  $O(g^4)$.
\item Construction of the continuum limit (OS reconstruction, 
  uniqueness).
\end{enumerate}

\subsection{Extension to $\mathrm{SU}(N)$}

At leading order, $\mathrm{SU}(N)$ contributes $(N^2 - 1)$ independent 
color channels, each behaving as $\mathrm{U}(1)$.  The non-abelian 
corrections---Baker--Campbell--Hausdorff commutators and metric 
curvature---enter at $O(g^4)$, the same order as the quartic.  
They modify the coefficient $C_0$ but not the scaling $W_0 = O(g^4)$, 
thanks to the antisymmetric cancellation of the $O(g^3)$ Christoffel term.

Numerical verification for $\mathrm{SU}(2)$ on the $4^3$ periodic lattice 
with $L = 2$ blocking confirms:
\begin{enumerate}
\item The initial entropy scales as $W_0^{\mathrm{SU}(2)} = 3 \times 
  W_0^{\mathrm{U}(1)}(\beta_{\mathrm{eff}})$ where 
  $\beta_{\mathrm{eff}} = \beta/4$.  At $\beta = 200$: 
  $W_0 \cdot \beta^2 = 197.6$, matching the prediction 
  $3 \times 4.05 \times 16 = 194.4$ to within $2\%$.

\item The dissipation fraction $D/W = 0.693 \pm 0.003$, identical to the 
  $\mathrm{U}(1)$ value, confirming that each color channel dissipates 
  independently at leading order.

\item The Baker--Campbell--Hausdorff correction to $W_0$ is $1.5\%$ at 
  $\beta = 100$ ($g^2 = 0.04$), consistent with an $O(g^6)$ correction 
  to an $O(g^4)$ quantity.
\end{enumerate}

For general $\mathrm{SU}(N)$ with the standard Wilson action 
$\beta_{\mathrm{lat}} = 2N/g^2$:
\[
  W_0^{\mathrm{SU}(N)} = 4N^2(N^2 - 1)\,C_1(d) / \beta_{\mathrm{lat}}^2.
\]

\subsection{What the framework achieves}

The entropy-controlled RG provides two categories of results:

\textbf{Structural results} (valid at all couplings):
\begin{enumerate}
\item Schur factorization $S = B^\top B$ (Proposition~\ref{prop:schur_factor}).
\item Exact comparability $R\,S\,R^\top = L^{-d}H_c$ 
  (Corollary~\ref{cor:exact_comp}).
\item Face-averaging cancellation $H_{yx}R^\top = 0$ 
  (Lemma~\ref{lem:face_cancel}).
\item Per-plaquette variance bound $P_{pp} \leq 1$ 
  (Theorem~\ref{thm:plaq_bound}).
\item KL chain rule: $W_{k+1} = W_k - D_k$ (Theorem~\ref{thm:KL}).
\end{enumerate}

\textbf{All-scale quantitative bounds} (for $\beta$ sufficiently large):
\begin{enumerate}\setcounter{enumi}{5}
\item Initial condition: $W_0 = C_0/\beta^2$ (Theorem~\ref{thm:W0}).
\item Dissipation bound at all scales: $D_k/W_k \geq 0.527$ 
  (Theorem~\ref{thm:diss}).
\item Geometric entropy decay: 
  $W_k \leq (1-f)^k C_0/\beta^2$ for all $k$ (Corollary~\ref{cor:gap}).
\item Exponential smallness at the AF scale: 
  $W_{k_{\max}} = O(e^{-c\beta}/\beta^2)$.
\item Cluster convergence: non-Gaussian corrections to correlation 
  functions decay exponentially at every scale, with 
  $\xi_{\mathrm{corr}} = O(1)$ in coarse lattice units 
  (Theorem~\ref{thm:cluster}).
\end{enumerate}

\textbf{Open:}
\begin{enumerate}\setcounter{enumi}{10}
\item \textbf{From entropy decay to mass gap.}  
  Section~\ref{sec:luyau} outlines a two-stage argument.
  \emph{Stage~1 (weak coupling, $k < k_*$):} the per-plaquette bound 
  controls the Holley--Stroock perturbation of the Gaussian 
  conditional, giving $\rho_k = \Omega(\beta_{\mathrm{eff}})$.  
  This is a \emph{single-block} property: the Gaussian conditional 
  variance is independent of the boundary condition, so no spatial 
  Dobrushin condition is needed.  
  \emph{Stage~2 (strong coupling, $k = k_*$):} at 
  $\beta_{\mathrm{eff}} \leq \beta_c$, the Wilson measure is a small 
  perturbation of Haar, and the Hessian Dobrushin 
  condition~\eqref{eq:dobrushin} holds; the 
  Martinelli--Olivieri theorem gives a volume-uniform spectral gap 
  $\Omega(1)$ in coarse units.  The resulting mass gap bound is 
  $m \cdot a \geq c/\beta^{1/d}$, a \emph{polynomial} lower bound 
  that is dramatically stronger than the earlier exponential estimate 
  $e^{-\beta/(2d)}$.  The remaining verification is the explicit 
  computation of~$\beta_c$ and the Hessian Dobrushin constants 
  for the block geometry.
\item Extension to infinite volume and $\mathrm{SU}(N)$.
\item Continuum limit (OS reconstruction).
\end{enumerate}

The mass gap problem has been reduced from a multi-scale \emph{control} 
problem (bounding the effective action at every scale) to a 
\emph{functional inequality} problem: the RG carries the theory from 
weak to strong coupling, with the per-plaquette bound ensuring the 
conditional gaps are controlled at every intermediate scale.  Spatial 
composition (Lu--Yau) is needed \emph{only} at strong coupling, 
where the Hessian Dobrushin condition holds by the smallness of 
$\beta_{\mathrm{eff}}$.  No step requires new techniques beyond the 
standard functional inequality 
literature~\cite{BGL,lu_yau,martinelli_olivieri}.

\subsection{The role of the Schur complement}

The iterated Schur complement $S^{(k)}$ serves as the 
\emph{reference geometry} for the RG flow, analogous to Perelman's 
Gaussian shrinker.  Its key properties are:
\begin{itemize}
\item It is explicitly computable (linear algebra on finite matrices).
\item It has a spectral fixed point: 
  $\lambda_{\min}(S^{(k)}_{\mathrm{phys}}) = 1$ for $L \geq 3$.
\item It provides the Gaussian gap that the true gap tracks.
\end{itemize}

The Gaussian shadow does not degrade under iteration---unlike 
Balaban's effective actions, which accumulate corrections that 
must be bounded term by term.  This is the main technical advantage 
of the entropy framework.

%% ================================================================
\section*{Acknowledgments}
%% ================================================================

I thank Claude (Anthropic) for extensive computational assistance 
in developing the numerical evidence presented in this paper.

\input{appendix_details}

\begin{thebibliography}{99}

\bibitem{JW}
A.~Jaffe and E.~Witten,
\emph{Quantum Yang--Mills theory},
Clay Mathematics Institute Millennium Prize Problem description, 2000.

\bibitem{Balaban1}
T.~Balaban,
\emph{Ultraviolet stability of three-dimensional lattice pure gauge field theories},
Comm.\ Math.\ Phys.\ \textbf{102} (1985), 255--275.

\bibitem{Balaban2}
T.~Balaban,
\emph{Renormalization group approach to lattice gauge field theories.
I.\ Generation of effective actions in a small field approximation and a coupling constant renormalization in four dimensions},
Comm.\ Math.\ Phys.\ \textbf{109} (1987), 249--301.

\bibitem{Balaban3}
T.~Balaban,
\emph{Large field renormalization. II.\ Localization, exponentiation, and bounds for the $\mathrm{R}$ operation},
Comm.\ Math.\ Phys.\ \textbf{122} (1989), 355--392.

\bibitem{DS1}
R.L.~Dobrushin and S.B.~Shlosman,
\emph{Completely analytical interactions: constructive description},
J.\ Stat.\ Phys.\ \textbf{46} (1987), 983--1014.

\bibitem{OS}
K.~Osterwalder and E.~Seiler,
\emph{Gauge field theories on a lattice},
Ann.\ Phys.\ \textbf{110} (1978), 440--471.

\bibitem{Fontes_antisym}
R.~da Costa Fontes,
\emph{Antisymmetric Christoffel cancellation in lattice Yang--Mills:
Schur complement positivity and block Poincar\'e inequalities},
Preprint, 2026.

\bibitem{gopfert_mack}
M.~G\"opfert and G.~Mack,
\emph{Proof of confinement of static quarks in $3$-dimensional $\mathrm{U}(1)$ 
lattice gauge theory for all values of the coupling constant},
Comm.\ Math.\ Phys.\ \textbf{82} (1982), 545--606.

\bibitem{lu_yau}
S.~L.~Lu and H.-T.~Yau,
\emph{Spectral gap and logarithmic Sobolev inequality for Kawasaki 
and Glauber dynamics},
Comm.\ Math.\ Phys.\ \textbf{156} (1993), 399--433.

\bibitem{martinelli_olivieri}
F.~Martinelli and E.~Olivieri,
\emph{Approach to equilibrium of Glauber dynamics in the one phase 
region.~I.~The attractive case},
Comm.\ Math.\ Phys.\ \textbf{161} (1994), 447--486.

\bibitem{BGL}
D.~Bakry, I.~Gentil, and M.~Ledoux,
\emph{Analysis and Geometry of Markov Diffusion Operators},
Grundlehren Math.\ Wiss., vol.~348, Springer, 2014.

\bibitem{Fontes2026_BE}
R.~da Costa Fontes,
\emph{Bakry--\'Emery curvature of the fractional Laplacian 
via fractional Brownian covariance},
Preprint, 2026.

\end{thebibliography}

\end{document}
